\section{Testbed Implementations} \label{sec:testbeds:main}


With the goal of developing, testing and validating network-aware applications and services, research initatives rely on experimentation platforms, commonly referred to as testbeds. This experimentation facilities provide access to network resources through virtualized 5G infrastructure, allowing third parties, such as \acsp{SP}, \acsp{SME}, vertical stakeholders or even the research community to experiment with solutions tailored to specific industry requirements~\cite{ICT-41-2020}. Such platforms abstract the complexities of managing the underlying infrastructure, allowing participants to focus on the design, implementation and evaluation of their own applications, rather than being preoccupied by the provisioning and operational details of the network.

The following section will delve into the details of the testbeds utilized in the literature for validating network-aware systems. Subsection \ref{sec:testbeds:5GC} performs an overview of the various components that tipically constitute testbeds or testing environments, highlighting the diversity of components and setups present in the literature. Then, Subsection \ref{sec:testbeds:5GC:5GC}, examines these testbeds regarding its network exposure capabilities, including adherence to key telecom standardization initiatives, such as CAMARA, \acs{CAPIF} and \acs{TMF} Open \acsp{API}. Lastly, the combined information of the researched testbeds and validation environments, including both the configuration setup and adherence to telecom standardization initiatives is summarized in Table \ref{tab:testbeds:description}.


\subsection{Testbed Architecture and Configurations} \label{sec:testbeds:5GC}

The analyzed testbeds follow a common architectural approach for supporting programmability and exposure of network capabilities, with their core modules being summarized in Table \ref{tab:testbeds:description}. At the center, a \acs{5G} core network forms the backbone of the \acs{5G} system, with implementations ranging from open-source projects, such as Open5GS, Free5GC, \acs{OAI} to commercial solutions, such as Amarisoft, Nokia, and Huawei. Then, the core network connects to the \acs{RAN}, via the transport network, for access to the end devices through \acs{gNB} base stations that may be deployed either as a physical component (e.g. \acs{OAI} \acs{RAN}, Ericsson and Huawei gNodeBs) or simulated software environments, such as UERANSIM.

To promote flexibility and reproducible results, testbeds heavily rely on virtualization and orchestration technologies, with most testbeds leveraging containerization platforms like Kubernetes. \acf{VM} environments are also common utilizing solutions, constructed via OpenStack and VirtualBox.  Moreover, \acs{MANO} solutions compliant with \acs{ETSI} \acs{NFV} standards are also present to manage network functions and services, with \acf{OSM} being widely utilized.


Regarding network capability exposure, an analysis was performed covering the various network exposure mechanisms across the various testbeds, namely the \acs{NEF} implementations discussed in Section \ref{exposure:NEFimplementations}, as well as adherence to the telecom standardization initiatives mentioned in Section \ref{sota:telcoinitiatives} (i.e. CAMARA, \acs{CAPIF}, \acs{TMF} Open \acsp{API}). A detailed analysis of the exposure mechanisms constructed utilizing this interfaces is provided in the following section.

\hide
{

Regarding physical implementations, multiple testbeds utilize commercial solutions for the \acs{RAN}. \citet{OnDemandNetworkQoSAdaptation} deploy their work on the experimetal infrastructure of the Dutch National \acs{6G} testbed (FNS), which includes a \acs{B5G} core and an Ericsson \acs{gNB} with outdoor antennas. In Aveiro, several contributions rely on the 5GAIner testbed, an Huawei commercial-grade \acs{5G SA} Network (3GPP release 17), featuring Huawei \acs{5G} New Radio cells suitable for both outdoor and indoor scenarios~\cite{5GAIner}. Alternatively, enumerous testing environments opt to simulate both the \acs{UE} and the \acs{gNB} by using UERANSIM~\citep{LoadBalancingConnectedVehicleServices,CloudEnabledDeployment5GN,DesignAndEvaluationNetworkDigitalTwin,NSSFfunction6GNetworksMLOPS}, an open source \acs{5G} \acs{UE} and \acs{RAN} simulator. However, UERANSIM does not actually implement \acs{5G} radio protocols bellow the \acs{RRC} layer, simulating it with a \acs{UDP}-based \acf{RLS} protocol to connect the \acs{gNB} and \acs{UE}, which may introduce disparities in the testing process providing unrealistically high throughput and overly optimistic latency due to the absence of physical layer processing~\cite{EvaluatingOpenSource5GSATestbeds}.

For experiments requiring real \acs{RF} transmission and reception, full \acs{RAN} implementations such as \acs{OAI} can be used, as it supports physical layer operations and is compatible with \acs{USRP} hardware. \citet{ExplorationNWDAFDevelopmentArchitecture6G} combines a Open5GS core and an \acs{OAI} \acs{RAN} with \acs{USRP} hardware to simulate real-world scenarios with a \acs{gNB} and \acs{UE}. As the \acs{OAI} also provides a 3GPP compliant \acs{OAI} \acs{5G} core (\acs{SA}), other authors utilize a full \acs{OAI} \acs{5G} system, with both the \acs{RAN} and \acs{5GC}~\cite{BenefitsCaveatsQoDNetworkAPIsVideoStreaming}. 

Another approach, widely used in the literature, correspond to network emulators. In this field, commercial offerings by Amarisoft are the most prominent in the literature, with multiple versions of Amarisoft Callbox being employed for testing of the \acs{RAN} and/or the \acs{5G} \acs{CN}~\citep{AdvancingPredictiveSecurity,PerformanceComparisonNWDAFSecurity}. This may include \acs{5G SA} setups incorporated with \acs{SDR} elements and \acfp{RRH} for enhanced coverage~\cite{VirtualOnBoardUnitMigration}. Additionally, \acsp{UE} can also be emulated via Amarisoft's Simbox and connected to \acsp{gNB} from Amarisoft Callbox~\cite{DemoOnDemand5GEdgeVerticalsViaThirdPartyUPF}.

}


\subsection{Network Exposure Interfaces} \label{sec:testbeds:5GC:5GC}

The integration of network exposure interfaces within testbed implementations demonstrates varying degrees of architectural maturity and standardization compliance. This section expands on the configuration of the testbeds provided in Table \ref{tab:testbeds:description}, with particular emphasis on their configuration for exposing network capabilities. This may include the implementation of \acs{NEF}, CAMARA and \acs{TMF} \acsp{API}, as well as secure service exposure through \acs{CAPIF}.

\citet{DesignAndEvaluationNetworkDigitalTwin} implements a \acf{NDT} platform with Open5GS of a physical \acs{5G} testbed built with Amarisoft, enabling closed-loop automation and \acs{AI}/\acs{ML} driven network intelligence. With this goal, the authors repurpose the Katana Slice Manager, developing an \acs{NDT} orchestrator responsible for the provisioning, synchronization, and automation of \acs{NDT} instances. More specifically, Katana is integrated with \acs{OSM} to manage the deployment of Open5GS \acs{NSD}/\acs{KNF}, guaranteeing they are instantiated within the Kubernetes environment~\cite{videoNDTDemokritos}. The system utilizes Prometheus for managing the collection, storage, and retrieval of network telemetry and simulation data, with visual analytics being obtained via Grafana. Although the authors specify an \acs{NDT} framework suitable for \acs{B5G}/\acs{6G} networks, it is not addressed the exposure of services suitable for third parties, with no references to either \acs{NEF} or CAMARA \acsp{API}.

In contrast, \citet{BenefitsCaveatsQoDNetworkAPIsVideoStreaming} utilizes the CAMARA \acs{QoD} \acs{API} to improve the \acs{QoE} of video streaming scenarios in a network with challenging conditions. This is achieved through a cache server that interfaces with the CAMARA \acs{QoD} \acs{API} to request temporary bandwidth boosts. However, the underlying implementation relies on low-level xApps deployed at the \acs{RIC}, a layer above the \acs{RAN}, with no references to either \acs{NEF} or other \acsp{NF} of the \acs{5GC}. This approach promotes tighter coupling between the application and the \acs{RAN} specific implementation, not taking advantage of the exposure and abstraction layer provided by the \acs{NEF}.

There are other approaches for offerring CAMARA services, incorporating additional abstraction layers that decouple the service logic from the underlying RAN implementation, focusing for instance in \acs{TMF} \acsp{API}~\cite{CAMARAaaSdemo}. The introduced \acs{CAMARAaaS} approach integrates CAMARA with \acs{TMF} \acsp{API} to offer a standardized interface to manage \acs{5G} resources, assuming network services can be exposed through \acs{TMF} \acsp{API}. More specifically, it showcases how the \acs{QoS} of an \acs{UE} can be controlled with the CAMARA \acs{QoD} Provisioning \acs{API}. In practice, it introduces a middleware layer to translate the CAMARA \acs{API} requests into operations from the \acs{TMF} 638 - Service Inventory Management \acs{API}, included in the suit of \acs{TMF} \acsp{API} offered and fulfilled by the \acs{OSL} \acs{OSS}. 

\citet{DynamicQoSAdaptationviaExposedServiceAPIs} extends on this use case by detailing end-to-end specifications for the implementation of two CAMARA services, the Device Location and \acs{QoD} \acsp{API}, in the context of a vehicular remote monitoring digital twin use case. The CAMARA \acsp{API} are exposed by the \acs{B5G} core at the \acs{AF} layer, enabling device location notifications and on-demand \acs{QoS} requests for the selected \acsp{UE}. By combining these APIs, a \acs{V2X} server can automatically request a high-priority \acs{QoS} profile whenever the Device Location \acs{API} indicates a vehicle has entered a congested geographical area, solely through the \acs{AF} layer without directly interacting with underlying network functions. This layer converts the requests to internal 3GPP procedures required by the \acs{NEF}, which coordinates with the relevant \acsp{NF}, i.e. the \acs{PCF} and the \acs{AMF} for the \acs{QoD} and Device Location \acsp{API}, respectively. This is especially relevant as even lower-level interactions, between the \acs{AMF} and the \acs{gNB} for obtaining the actual \acs{NGAP} Location Reports, get abstracted.

With respect do dynamic \acs{QoS} adaptation via exposed \acsp{API}, \cite{NextGenConnectivityDynamicQoS} proposes two EdgeApps, the Situational Awareness EdgeApp and the Quality Awareness EdgeApp, responsible for requesting modifications in the \acs{QoS} (requesting a slice profile with higher bandwidth) upon detection of special events using geofencing (e.g. entering higher \acs{QoS} area) and antecipating potential network deterioration, respectively. It differs from the previous work, by introducing Nokia's \acf{NaC} solution to expose the CAMARA \acs{QoD} and Device Location \acsp{API}. The \acs{NaC} \acs{SDK} is integrated directly into the edge network applications to create real-time events and trigger changes in the \acs{QoS}, which are fulfilled by the \acs{NaC} backend services through the \acs{NEF} into the Telenet 5G Core, also based on Nokia technology. Nonetheless, the details of the underlying implementation at the level of the \acs{5G} core are not specified, with the solution being tighly coupled to the proprietary Nokia \acs{NaC}.

A common limitation across all the previous analysed implementations is the lack of \acs{CAPIF} support for secure discovery, exposure and consumption of northbound \acsp{API}, provided by \acs{NEF} and/or CAMARA. \cite{PathwayNaaSCAMARA} presents a solution that incorporates \acs{CAPIF} to publish CAMARA \acsp{API}, and policy their exposure to third party applications. It leverages the CAMARA \acs{QoD} \acs{API} to create a traffic management application to be used by third parties (e.g. streaming platform), modifying the \acs{QoS} profile (bandwidth and latency) associated to individual \acsp{UE} to prevent performance degradation. The \acs{QoS} modifications are abstracted by CAMARA and enforced by the \acs{NEF}, which interfaces with the \acs{5GC} via its \texttt{AsSessionWithQoS} \acs{API}. This corresponds to the most complete test environment, regarding the exposure of network capabilities, incorporating \acs{NEF}, CAMARA and \acs{CAPIF} in the final solution, as shown in Table \ref{tab:testbeds:description}.

\begin{table}[ht]
\centering
\tiny
\renewcommand{\arraystretch}{1.3} % Increase row height
\rowcolors{2}{gray!15}{white}     % Alternating row colors
\caption{Testbed Implementation Configurations}
\label{tab:testbeds:description}
\begin{tabular}{p{2cm} p{2cm} p{2cm} p{2cm} p{1.4cm} p{1.4cm} p{1.4cm} p{1.4cm}}
\toprule
\textbf{Source} & \textbf{5G Core} & \textbf{5G RAN} & \textbf{Virtualisation / Orchestration} & \textbf{NEF} & \textbf{CAMARA} & \textbf{CAPIF} & \textbf{TMF} \\
\midrule
Advancing Predictive Security~\cite{AdvancingPredictiveSecurity} & Amarisoft Classic, NCSRD-DS-5GDDoS testbed. & Amarisoft Callbox Mini and 2 cells - Amarisoft Classic. & Not specified. & NO. & NO. & NO. & NO. \\
Exploration of \acs{NWDAF} Development Architecture~\cite{ExplorationNWDAFDevelopmentArchitecture6G} & Open5GS. & \acs{OAI} using \acs{USRP} B210. & Kubernetes. & NO. & NO. & NO. & NO. \\
Network Slice Selector in \acs{IoT} Services~\cite{NetworkSliceSelectorIoT} & OpenAirInterface, GNS3 emulator. & \acs{OAI} gNodeB (CUPS Ubuntu). & Kubernetes (openair-k8s), OSM, GNS3. & NO (NWDAF focus). & NO. & NO. & NO. \\
Cloud-Enabled Deployment of \acs{5G} Core Network~\cite{CloudEnabledDeployment5GN} & Open5GS. & UERANSIM (\acs{gNB} and \acs{UE}). & VirtualBox, Debian \acsp{VM}. & Custom Quasi-NEF in Python. & NO. & NO. & NO. \\
Dynamic QoS Adaptation Via Exposed Service \acsp{API}~\cite{DynamicQoSAdaptationviaExposedServiceAPIs} & TNO B5G Core (Open5GS based), TNO Hague lab, Rel-17. & Ericsson gNodeB, \acs{SA} mode. & OpenStack \acsp{VM}. & YES (TNO NEF). & \acs{QoD}, Device Location. & NO. & NO. \\
Load Balancing in Connected Vehicle Services~\cite{LoadBalancingConnectedVehicleServices} & Free5GC, distributed edge servers, v3.3.0. & UERANSIM (\acs{UE} and \acs{gNB}). & Ubuntu VMs. & Custom Go implementation. & Traffic Influence \acs{API}. & NO. & NO. \\
An Automated Testing and Orchestration Framework~\cite{AutomatedTestingOrchestrationFramework} & Huawei \acs{5G} \acs{SA} Core, 5GAIner infrastructure, Rel-17. & Huawei 5G New Radio cells. & \acs{OSM}, Kubernetes, OpenStack. & NEFSim~\cite{NEFSimMediaNetLabGithub}. & NO. & NO. & YES (\acs{OSL}). \\
Virtual On-Board Unit Migration~\cite{VirtualOnBoardUnitMigration} & GAIA 5G (Amarisoft), University of Murcia. & Amarisoft macro cells. & OpenStack (Rocky), \acs{OSM}. & NEFSim~\cite{NEFSimMediaNetLabGithub}. & NO. & NO. & NO. \\
Next-Gen Connectivity: Dynamic \acs{QoS} Optimization~\cite{NextGenConnectivityDynamicQoS} & Telenet \acs{5G} \acs{SA} (Nokia), Antwerp VITAL-5G testbed. & Antwerp gNodeB units. & OpenStack, OSM (VITAL-5G). & YES (with Nokia \acs{NaC}). & \acs{QoD}, Device Location. & NO. & NO. \\
On the benefits and caveats of exploiting \acs{QoD}~\cite{BenefitsCaveatsQoDNetworkAPIsVideoStreaming} & OpenAirInterface \acs{5G} Core. & OpenAirInterface \acs{5G} \acs{RAN}. & Not specified. & NO (uses \acs{RIC} directly). & \acs{QoD}. & NO. & NO. \\
On-Demand Network QoS Adaptation~\cite{OnDemandNetworkQoSAdaptation} & TNO B5G Core, Dutch National 6G testbed (FNS), Rel-17. & Ericsson gNodeB, Rel-16. & TNO research cloud. & YES (Internal TNO B5G). & \acs{QoD}, Performance Metrics. & NO. & NO. \\
Design and Evaluation of a \acs{NDT}~\cite{DesignAndEvaluationNetworkDigitalTwin} & Open5GS, NCSR Demokritos \acs{5G}+ testbed. & UERANSIM. & \acs{NDT} Katana Orchestrator. & NO. & NO. & NO. & NO. \\
\acs{NSSF} Function in \acs{6G} Networks Based on MLOps~\cite{NSSFfunction6GNetworksMLOPS} & Not specied, containerised. & UERANSIM simulator. & Kubernetes (K8s cluster). & NO (NSSF/NWDAF focus). & NO. & NO. & NO. \\
Performance Comparison of NWDAF-Based Security~\cite{PerformanceComparisonNWDAFSecurity} & Amarisoft Callbox Classic, "\acs{5G}-in-a-box" lab platform. & Amarisoft (\acs{SA} mode). & NO. & NO. & NO. & NO. & NO. \\
Offering CAMARA APIs as a Service - Demo~\cite{CAMARAaaSdemo} & Huawei \acs{5G} Core, 5GAIner testbed, Rel-17. & 5GAIner New Radio cells. & Kubernetes (\acs{OSL}'s CRIDGE). & NO. & \acs{QoD} Provisioning \acs{API}. & NO. & YES (TMF 638, \acs{OSL}). \\
Pathways towards network-as-a-service~\cite{PathwayNaaSCAMARA} & \acs{SA}, Release 15. & single gNodeB. & containerized. & YES. & \acs{QoD} \acs{API}. & YES. & NO. \\
Experiences and challenges on testing~\cite{ExperiencesChallengesTestingValidatingNetworkAware} & Patras5G Core, Patras5G testbed. & AW2S gNB. & containerized and virtualized. & NO. & NO. & NO. & YES (\acs{OSL}). \\
Pushing the Boundaries of Scalable 5G Core~\cite{pushingBoundariesScalable5GCN} & Open5GS. & Not specified. & Kubernetes cluster. & custom microservices \acs{NEF}. & NO. & YES (OpenCAPIF). & NO. \\
Optimizing \acs{3GPP} \acs{5G} \acs{NEF} Integration~\cite{optimizing5GNEFIntegrationThroughNEFandUPF} & Not specified. & General Access Nodes. & Not specified. & YES (custom \acs{NEF}). & \acs{QoD} API. & NO. & NO. \\
Capability Exposure Vitalizes \acs{5G}~\cite{CapabilityExposureVitalizes5GNetwork} & 5G Core Trial, Hangzhou site. & 5G RAN. & Not specified. & YES (custom). & NO. & NO. & NO. \\
Data Collection Using \acs{NWDAF}~\cite{DataCollectionUsingNWDAF5GCN} & 5G Core, Linux-based testbed. & Not specified. & Linux \acs{VM} and Namespaces. & NO. & NO. & NO. & NO. \\
\acs{TMF} Open Digital Architecture for \acs{5G} Service Orchest.~\cite{TMFOpenDigitalArchitecture5GServiceOrchestration} & Huawei \acs{5G} \acs{SA} Core, 5GAIner, Rel-17. & Huawei \acs{5G} New Radio cells. & \acs{OSM}, Kubernetes, OpenStack. & NO. & \acs{QoD} Provisioning \acs{API}. & NO. & YES (\acs{OSL}). \\
\bottomrule
\end{tabular}
\end{table}


